\chapter{Cài đặt PODEM}
\label{sec:p5-cai-dat-podem}
\textbf{Kiến trúc.}
Lõi chương trình hiện thực PODEM~\cite{p3_goel1981}, nhận \texttt{Circuit},
\texttt{Fault}, giới hạn quay lui \texttt{max\_backtracks} (mặc định 1000)
và tùy chọn trace. Sau mỗi lần gán PI, \texttt{imply} khởi tạo PI chưa gán
bằng $X$, cấy lỗi và mô phỏng năm giá trị $0,1,X,D,D'$ theo thứ tự topo;
pattern PI chỉ chứa $0,1,X$. \texttt{unroll}
(Chương~\ref{sec:p4-atpg-tuan-tu}) biến mạch tuần tự thành mạch tổ hợp
nhiều khung; \texttt{run} đọc netlist, gọi PODEM và kiểm tra kết quả bằng
bộ mô phỏng lỗi \texttt{fault\_sim}. Cổng và đầu vào được chọn theo thứ tự
netlist ổn định, chưa dùng SCOAP.

\textbf{Objective và backtrace.}
Trước kích hoạt, objective yêu cầu net lỗi bằng giá trị ngược stuck-at.
Sau kích hoạt, PODEM chọn cổng D-frontier đầu tiên còn X-path tới PO và
input $X$ đầu tiên của nó; AND/NAND cần 1, OR/NOR cần 0 để truyền sai biệt.
Backtrace đưa mục tiêu về PI: NAND/NOR/NOT đảo giá trị, AND/OR/BUFF giữ
nguyên. Với XOR/XNOR, objective lan truyền chọn input $X$ đầu tiên và thử 0;
khi backtrace gặp cổng còn đúng một $X$, giá trị được tính từ parity rail tốt
đã biết ($D\mapsto1$, $D'\mapsto0$; $q=0$ với XOR, $q=1$ với XNOR). Còn nhiều
$X$ thì chọn $X$ đầu tiên, thử 0 và để imply tính lại, không giả định các
$X$ còn lại bằng 0.

\textbf{Quay lui, trace và điều kiện dừng.}
Khi nhánh thất bại, decision stack bỏ các quyết định con, đảo PI gần nhất
còn giá trị chưa thử rồi imply lại toàn mạch. Trace xuất một hàng bảy cột
cho mỗi lần gán hoặc đảo PI; hàng dẫn tới ngõ cụt ghi \texttt{backtrack},
hàng đảo quyết định dùng ``—'' cho objective/backtrace, khớp quy ước golden
trace ở Chương~\ref{sec:p3-podem}. Trace thật trên c17 trùng hai bước ở
Bảng~\ref{tab:p3-c17}; chương trình có công cụ xuất trace đầy đủ.
PODEM trả \texttt{DETECTED} khi PO có $D/D'$; sau kích hoạt, nhánh không còn
D-frontier có X-path bị loại. Hết cây quyết định là \texttt{UNTESTABLE};
chạm \texttt{max\_backtracks} là \texttt{ABORTED}, không phải
\texttt{UNTESTABLE}, vì chưa thử hết mọi vector.

\textbf{Cấy lỗi và fault-frame.}
Stem tác động mọi fanout; branch chỉ tác động ngõ vào cổng đích.
Cấy lỗi giữ rail tốt, ép rail lỗi: $D+SA0=D$, $D+SA1=1$,
$D'+SA0=0$, $D'+SA1=D'$. Trên mạch trải, imply cấy đồng thời danh sách
fault-instance; objective/backtrace lần lượt chọn fault-frame làm mục tiêu.
PODEM có thể gán $Q@0$, nên bộ kiểm chứng tuần tự phải phân biệt kết quả
có điều kiện và chuỗi bảo đảm, không tự coi trạng thái đầu là điều khiển được.

\textbf{Kiểm thử và giới hạn.}
Logic năm giá trị và xử lý lỗi theo mô hình trong giáo trình
\cite[Ch.~3--4]{p1_wang2006}. Input sai (net/nhánh không thuộc mạch,
$SV\notin\{0,1\}$, giới hạn âm) bị từ chối bằng \texttt{ValueError}.
Bộ kiểm thử đạt 350 test (Python 3.12.15, pytest 9.1.1), riêng logic và
PODEM 255 test. Với c17 $11/SA0$: $X10XX$, 0 backtrack, hai phép gán PI;
mạch phụ $t/SA0$: $01$, một backtrack (giới hạn 0 trả \texttt{ABORTED}).
Bộ mô phỏng lỗi độc lập xác nhận các mẫu. Chưa hỗ trợ lỗi nhánh qua DFF, FAN hay
heuristic SCOAP.
